\section{From Hilbert vectors to the original smooth automorphic map}
\label{sec:smooth-representatives}

This chapter proves the original smooth-representative theorem and the
full \texttt{BackForm.\allowbreak exists\_\allowbreak multMap} on its stated inputs.  The
actual adelic chart and fixed-order growth estimate are supplied in
ASB; the representative construction uses only its local chart
comparison, not the global rational-point count.
The final installation retains the original initial boundedness and
target finite-volume inputs.  The common datum \(R\) is never replaced.

We first work with a locally compact, second-countable Hausdorff group
\(G\), a left Haar measure \(\mu\), a complex Hilbert space
\(\mathcal H\), and a representative assignment
\(\iota:\mathcal H\to\mathbf C^G\).
The two assumptions on this assignment are
\[
\begin{split}
 &\iota(av+w)=a\,\iota(v)+\iota(w)\quad\mu\text{-almost everywhere},\\
 &\forall K\subset G\text{ compact}\ \exists M_K\ge0\
   \forall v\in\mathcal H:\quad
   \iota(v)\in L^1(K),\qquad
   \int_K|\iota(v)|\,d\mu\le M_K\|v\|.
\end{split}
\tag{SPR.1}
\]
An equality between representatives below means a globally Haar
almost-everywhere equality, not an equality restricted to one
fundamental domain.  Write
\(\operatorname{Rep}(v,\phi)\) for
\(\iota(v)=\phi\) almost everywhere.
When a representation is used, it is a homomorphism
\(\rho:G\to\operatorname{Isom}_{\mathbf C}(\mathcal H)\) with
strongly continuous orbits and
\[
 \operatorname{Rep}\bigl(\rho(g)v,\iota(v)(\,\cdot\,g)\bigr).
 \tag{SPR.2}
\]
The group is assumed unimodular precisely where right translation of
almost-everywhere identities is used.  These assumptions hold on the
original quotient by RQC1--RQC6 and RQC11.

% BEGIN SPR1
\paragraph{SPR1: local \(L^1\) calculus and uniqueness of continuous representatives.}
\label{claim:SPR1}
Under (SPR.1), the map
\[
 I_K:\mathcal H\longrightarrow L^1(K,\mu),\qquad
               v\longmapsto[\iota(v)|_K]
 \tag{SPR.3}
\]
is complex linear and bounded, with norm at most \(M_K\).
It commutes with every Bochner integral for which the integrand is
integrable in \(\mathcal H\).  Continuous representatives are unique.

\emph{Proof.}
The first line of (SPR.1) implies the zero, addition and scalar
identities in the quotient \(L^1(K)\); the second gives its norm
bound.  Approximate a Bochner-integrable function by integrable simple
functions.  Linearity gives commutation for the simple functions,
and boundedness of \(I_K\) passes the identity to their \(L^1\) limits.

For completeness, the global representative test used below is
\[
 \left[\ \int_K\iota(v)\,d\mu=\int_K\phi\,d\mu
           \quad\hbox{for every compact }K\ \right]
 \quad\Longrightarrow\quad \operatorname{Rep}(v,\phi),
 \tag{SPR.4}
\]
provided \(\phi\) is locally integrable.
Indeed the locally integrable difference defines a complex measure
on each relatively compact open set.  Haar regularity and
integrability of its density approximate every measurable set there
from compact sets in the total-variation measure.  All its compact
integrals being zero, this complex measure is zero.  Apply this
separately to the positive and negative parts of the real and
imaginary densities, and cover \(G\) by countably many relatively
compact sets.  The difference is zero almost everywhere on \(G\).

If two continuous representatives agree almost everywhere, their
difference vanishes everywhere.  A nonzero value would, by
continuity, give a nonempty open set on which its absolute value is
bounded below by a positive constant.  Every nonempty open set has
positive Haar measure.  This proves uniqueness, including when
\(\iota\) is not injective as a map of vectors.
These arguments give the entire local representative calculus used
by \texttt{LocRep.\allowbreak exists\_\allowbreak pairing} and \texttt{rep\_\allowbreak of\_\allowbreak pairing}.
% END SPR1

% BEGIN SPR2
\paragraph{SPR2: a continuous parameter integral represents the Hilbert integral.}
\label{claim:SPR2}
Let \(Q\) be a compact subset of a second-countable Hausdorff parameter
space with locally finite Borel measure \(\nu\).  Suppose \(w(t)\) is
continuous with values in \(\mathcal H\), the function
\((t,x)\mapsto\phi(t,x)\) is continuous, and
\(\operatorname{Rep}(w(t),\phi(t,\cdot))\) holds for each \(t\).
Then
\[
 \Phi(x)=\int_Q\phi(t,x)\,d\nu(t)
 \quad\hbox{is continuous},\qquad
 \operatorname{Rep}\left(\int_Qw(t)\,d\nu(t),\Phi\right).
 \tag{SPR.5}
\]

\emph{Proof.}
Near a fixed \(x_0\), choose a compact neighborhood \(B\).
The continuous integrand is bounded on \(Q\times B\), so dominated
convergence proves continuity in \(x\).
For any compact \(K\), both \(w\) and the scalar integrand on
\(Q\times K\) are integrable.  The bounded functional
\[
 v\longmapsto\int_K\iota(v)\,d\mu
\]
from SPR1 commutes with the Hilbert integral.  Fubini and the
individual representative identities yield
\[
 \int_K\iota\!\left(\int_Q w(t)\,d\nu(t)\right)d\mu
 =\int_Q\int_K\phi(t,x)\,d\mu(x)\,d\nu(t)
 =\int_K\Phi(x)\,d\mu(x).
\]
SPR1's compact-set test proves (SPR.5).
No common exceptional set for all parameters was assumed:
integration is performed through the continuous functionals.
This proves the full original
\texttt{representative\_\allowbreak parametric\_\allowbreak integral\_\allowbreak continuous\_\allowbreak family}.
% END SPR2

% BEGIN SPR3
\paragraph{SPR3: compact averaging produces a continuous representative.}
\label{claim:SPR3}
Assume (SPR.1)--(SPR.2).  For \(\eta\in C_c(G)\) and \(v\in\mathcal H\),
the vector
\[
 A_\eta v=\int_G\eta(y)\rho(y)v\,d\mu(y)
\]
has the continuous representative
\[
 \phi_{\eta,v}(x)
   =\int_G\eta(y)\iota(v)(xy)\,d\mu(y)
   =\int_G\eta(x^{-1}z)\iota(v)(z)\,d\mu(z).
 \tag{SPR.6}
\]
This conclusion needs left Haar invariance only.

\emph{Proof.}
The Hilbert integral converges since its integrand is continuous,
compactly supported, and has norm \(|\eta(y)|\|v\|\).
The second expression follows by the left translation \(z=xy\).
For \(x\) in a compact neighborhood \(B\) of \(x_0\), its scalar
integrand is supported in \(B\operatorname{supp}\eta\) and dominated
by
\(\|\eta\|_\infty|\iota(v)(z)|\) there.
This is integrable by (SPR.1).  Dominated convergence gives
continuity of \(\phi_{\eta,v}\).

Fix a compact \(K\).  From (SPR.2), the local bound and isometry,
\[
 \int_K|\iota(v)(xy)|\,d\mu(x)
 =\int_K|\iota(\rho(y)v)(x)|\,d\mu(x)
 \le M_K\|v\|.
\]
Consequently the absolute double integral of
\(\eta(y)\iota(v)(xy)\) over \(G\times K\) is at most
\(M_K\|v\|\|\eta\|_1\).
The product function is measurable up to completion: choose a
measurable representative of \(\iota(v)\), and use the Haar
null-set property for the multiplication map, justified slice by
slice by left translation and Fubini.
Thus Fubini is applicable.  The same compact integration
functional as in SPR2 gives
\[
 \int_K\iota(A_\eta v)\,d\mu=\int_K\phi_{\eta,v}\,d\mu.
\]
SPR1 proves the representative identity.
This is the original
\texttt{exists\_\allowbreak continuous\_\allowbreak representative\_\allowbreak average}, with no
smoothness assumption on \(v\).
% END SPR3

% BEGIN SPR4
\paragraph{SPR4: continuous representatives are dense at the same fixed finite level.}
\label{claim:SPR4}
Let \(H,U\le G\), with \(U\) open, and suppose
\[
          y^{-1}Hy\subseteq H\qquad(y\in U).
 \tag{SPR.7}
\]
Every \(H\)-fixed vector is a norm limit of \(H\)-fixed vectors
having continuous representatives.

\emph{Proof.}
For an \(H\)-fixed \(v\) and \(\epsilon>0\), strong continuity gives
an identity neighborhood
\[
 W=U\cap\{y:\|\rho(y)v-v\|<\epsilon/2\}.
\]
Choose a nonnegative continuous compactly supported function
\(f\) supported in \(W\), with \(f(1)=1\).
Such a function follows from local compactness and the separation
of a compact neighborhood of \(1\) from the complement of \(W\);
the standard separation construction gives a function between
zero and one.  Its integral \(c\) is positive by Haar positivity
and finite by compact support.
Set \(\eta=f/c\).  Then \(\int\eta=1\), and SPR3 gives a
continuous representative for \(A_\eta v\).

For \(h\in H\) and \(y\) in the support,
\[
 \rho(h)\rho(y)v
 =\rho(y)\rho(y^{-1}hy)v=\rho(y)v
\]
by (SPR.7).  Hence \(A_\eta v\) is fixed by the same \(H\).
Moreover
\[
 \|v-A_\eta v\|
 \le\int_G\eta(y)\|v-\rho(y)v\|\,d\mu(y)
 \le\epsilon/2<\epsilon.
\]
Choosing \(\epsilon\to0\) proves density.
This checks the original
\texttt{exists\_\allowbreak fixed\_\allowbreak average\_\allowbreak approx\_\allowbreak of\_\allowbreak subgroup\_\allowbreak fixed}
with its actual normalization condition; no smaller varying
finite level has been introduced.
% END SPR4

% BEGIN SPR5
\paragraph{SPR5: bounded local evaluations extend to the Hilbert closure.}
\label{claim:SPR5}
Let \(A:\mathcal H\to\mathcal H\) be bounded complex linear and
\(E\subseteq\mathcal H\) a complex subspace.
Suppose every \(Av\), \(v\in E\), has a continuous representative
\(\psi_v\), and for every compact \(B\subset G\) there is \(c_B\ge0\)
such that
\[
           |\psi_v(x)|\le c_B\|v\|
             \qquad(v\in E,\ x\in B).
 \tag{SPR.8}
\]
Then there are bounded linear functionals
\(\Lambda_x:\mathcal H\to\mathbf C\) such that, for \(v\in\overline E\),
\[
 x\longmapsto\Lambda_xv\text{ is continuous},\qquad
             \operatorname{Rep}(Av,\Lambda_\bullet v).
 \tag{SPR.9}
\]

\emph{Proof.}
SPR1 and (SPR.1) make \(v\mapsto\psi_v(x)\) linear:
both sides of an addition or scalar identity are continuous
representatives of the same Hilbert vector.
The singleton case of (SPR.8) makes this functional bounded on
\(E\); extend it by continuity to \(\overline E\).
Compose with the orthogonal projection
\(P:\mathcal H\to\overline E\) to obtain \(\Lambda_x\) on all of
\(\mathcal H\), with the same bound.
The needed projection is obtained directly by minimizing distance
to the closed subspace: a minimizing sequence is Cauchy by the
parallelogram identity; its limit exists by completeness, and
differentiation of its squared distance along the subspace gives
orthogonality and the norm bound \(\|P\|\le1\).

For \(v_n\in E\) with \(v_n\to v\in\overline E\),
\[
 \sup_{x\in B}|\Lambda_xv_n-\Lambda_xv|
                       \le c_B\|v_n-v\|.
 \tag{SPR.10}
\]
Thus \(\Lambda_\bullet v\) is a locally uniform limit of continuous
functions and is continuous.  Also \(Av_n\to Av\), and the bounded
map \(I_K\) of SPR1 gives
\(\iota(Av_n)\to\iota(Av)\) in \(L^1(K)\).
Equation (SPR.10) gives convergence of their continuous
representatives in the same \(L^1(K)\).
Their two limits agree almost everywhere on \(K\), and a countable
compact cover gives (SPR.9) globally.
This proves the full original \texttt{exists\_\allowbreak eval\_\allowbreak closure}.
Only the evaluations on \(\overline E\) matter; operator-norm
continuity of \(x\mapsto\Lambda_x\) is not required.
% END SPR5

% BEGIN SPR6
\paragraph{SPR6: evaluation of a compact parameter average at the original level.}
\label{claim:SPR6}
Assume also that \(\mu\) is right invariant.
Let \(H,U\) satisfy (SPR.7), with \(U\) open.
Let \(Q\subset\mathbf R^N\) and \(\Omega\subset G\) be compact,
\(\Theta:\mathbf R^N\to G\) continuous, and \(C\ge0\).
Suppose, for every continuous nonnegative right-\(H\)-invariant
function \(f\),
\[
             \int_Q f(\Theta(t))\,dt
                  \le C\int_\Omega f(y)\,d\mu(y).
 \tag{SPR.11}
\]
For every continuous scalar \(\kappa\) there is a family of bounded
linear evaluations \(\Lambda_x\) such that every \(H\)-fixed \(v\)
satisfies
\[
 \operatorname{Rep}\left(
       \int_Q\kappa(t)\rho(\Theta(t))v\,dt,\,
       x\longmapsto\Lambda_xv\right),
 \qquad \Lambda_\bullet v\text{ is continuous}.
 \tag{SPR.12}
\]

\emph{Proof.}
The Hilbert averaging operator \(A\) is bounded by
\(\int_Q|\kappa(t)|\,dt\).
Let \(E\) be the subspace of \(H\)-fixed vectors having continuous
representatives.  Its linearity follows from SPR1.
For \(v\in E\), write \(\operatorname{Rep}(v,\phi)\).
Right invariance of Haar measure, (SPR.2), and continuous
uniqueness give
\[
                         \phi(xh)=\phi(x)\qquad(h\in H).
 \tag{SPR.13}
\]
SPR2 identifies the continuous representative of \(Av\) as
\(\psi_v(x)=\int_Q\kappa(t)\phi(x\Theta(t))\,dt\).

For \(x\) in a compact \(B\), apply (SPR.11) to
\(f(y)=|\phi(xy)|\).  Put
\(b=\sup_{t\in Q}|\kappa(t)|\).  Left Haar invariance and (SPR.1)
give
\[
\begin{split}
 |\psi_v(x)|
 &\le bC\int_\Omega|\phi(xy)|\,d\mu(y)\\
 &\le bC\int_{B\Omega}|\phi(z)|\,d\mu(z)
 \le bC M_{B\Omega}\|v\|.
\end{split}
\tag{SPR.14}
\]
The second inequality follows since \(x\Omega\subseteq B\Omega\);
no invariance of \(\Omega\) is assumed.
Apply SPR5, then SPR4, since every \(H\)-fixed vector belongs to
\(\overline E\).  This is exactly the original
\texttt{SmoothRep.\allowbreak exists\_\allowbreak eval}.  In the adelic application ASB4
supplies (SPR.11) with the same prescribed finite level.
% END SPR6

% BEGIN SPR7
\paragraph{SPR7: the vector-valued one-dimensional averaging identity.}
\label{claim:SPR7}
Let \(c:\mathbf R\to G\) be a continuous one-parameter subgroup.
If
\(\frac{d}{dt}|_{t=0}\rho(c(t))w=w'\), then, for \(s>0\),
\[
 w=\int_0^s\frac1s\,\rho(c(r))w\,dr
       -\int_0^s\frac{s-r}{s}\,\rho(c(r))w'\,dr.
 \tag{SPR.15}
\]

\emph{Proof.}
The group law and boundedness of \(\rho(c(r))\) transport the
derivative at zero to
\[
 \frac{d}{dt}\bigg|_{t=r}\rho(c(t))w=\rho(c(r))w'.
\]
The latter is continuous in \(r\) by strong continuity.
Differentiate
\((s-r)\rho(c(r))w/s\), integrate from \(0\) to \(s\), and use
its endpoint values \(w\) and \(0\).
The Banach-valued fundamental theorem follows by integrating
the continuous derivative; it may also be checked after every
continuous real linear functional, which separates points.
All integrals are of continuous functions on a compact interval.
This proves the original weighted-average identity, with its
minus sign and exact scale.
% END SPR7

% BEGIN SPR8
\paragraph{SPR8: a finite cube identity preserving the noncommutative word order.}
\label{claim:SPR8}
Let \(c_1,\ldots,c_N\) be continuous one-parameter subgroups,
and let \(S\subseteq\mathcal H\) be stable under maps \(d_j\)
such that
\(\frac{d}{dt}|_0\rho(c_j(t))v=d_jv\) for every \(v\in S\).
Set
\[
 Q=[0,s]^N,\qquad
       \Theta(t)=c_1(t_1)\cdots c_N(t_N).
\]
For \(I=\{i_1<\cdots<i_r\}\subseteq\{1,\ldots,N\}\), put
\[
 \kappa_I(t)=
   \prod_{j\notin I}\frac1s
   \prod_{j\in I}\left(-\frac{s-t_j}{s}\right),
 \qquad D_Iv=d_{i_r}\cdots d_{i_1}v,\qquad D_\varnothing v=v.
\]
Then
\[
             v=\sum_{I\subseteq\{1,\ldots,N\}}
                  \int_Q\kappa_I(t)\rho(\Theta(t))D_Iv\,dt
                     \qquad(v\in S).
 \tag{SPR.16}
\]
In particular the coefficient functions and finite derivative
words are chosen before the maps \(d_j\), the stable set \(S\),
and the vector \(v\).

\emph{Proof.}
For \(N=0\), the cube is a point of zero-dimensional volume one,
and there is one empty word with coefficient one.
Assume the formula for \(N\).  In each term apply (SPR.15)
for \(c_{N+1}\) to the vector \(D_Iv\in S\).
The nondifferentiated term appends the factor \(1/s\).
The differentiated term appends the scalar
\(-(s-t_{N+1})/s\) and replaces \(D_Iv\) by
\(d_{N+1}D_Iv\).
The group action becomes
\(\rho(c_1(t_1)\cdots c_N(t_N)c_{N+1}(t_{N+1}))\).
Fubini for continuous Hilbert-valued functions on the compact cube
gives (SPR.16) at \(N+1\).

The new derivative is placed on the left of the old derivative
word; no derivative operators have been commuted.  In the source's
right-fold convention the corresponding list is
\([i_r,\ldots,i_1]\).
This explicit choice supplies every field of the original
\texttt{eq\_\allowbreak sum\_\allowbreak integral\_\allowbreak cube\_\allowbreak derivatives}.
% END SPR8


% BEGIN SPR9
\paragraph{SPR9: smooth curve products remain smooth after ordered differentiation.}
\label{claim:SPR9}
In this paragraph the global isometry and strong-continuity
assumptions on \(\rho\) are unnecessary: only the group law and
bounded linearity of each \(\rho(g)\) are used, together with
the following smoothness hypothesis.  No topology on \(G\)
is needed for this assertion.
Let \(A\subseteq P\) index maps \(e_p:\mathbf R\to G\), and put
\[
 \operatorname{Sm}(v)\ \Longleftrightarrow\
  \left[t\longmapsto
   \rho(e_{p_1}(t_1)\cdots e_{p_k}(t_k))v
     \text{ is }C^\infty
       \text{ for every finite admissible list }p_1,\ldots,p_k\right].
 \tag{SPR.17}
\]
Write \(d_pv=\frac{d}{dr}|_0\rho(e_p(r))v\) when
\(\operatorname{Sm}(v)\).
Then smooth vectors are stable under these derivatives and under
finite products of the actions \(e_p(r)\).  More precisely, for
every admissible word \(l\),
\[
 t\longmapsto d_l\!
       \left(\rho(e_{p_1}(t_1)\cdots e_{p_k}(t_k))v\right)
           \quad\hbox{is }C^\infty,
 \tag{SPR.18}
\]
where \(d_l\) has the original right-fold order.

\emph{Proof.}
Use the following strengthened smoothness property for a
parameter-dependent vector \(\Phi(z)\).
After every smooth change of its parameters and after multiplication
on the left by any finite product
\(\rho(e_{q_1}(a_1(s))\cdots e_{q_m}(a_m(s)))\), with admissible
\(q_j\) and smooth real functions \(a_j\), the resulting
Hilbert-valued function is smooth.
A constant vector \(v\) has this property by (SPR.17), composed
with the smooth coordinate map
\(s\mapsto(a_1(s),\ldots,a_m(s))\).
Adding a parameter-dependent finite group word preserves the
property by concatenation of the two words.

It is also preserved by \(\Phi\mapsto d_p\Phi\).
For a given outer word, introduce a new real variable \(r\) and
insert \(e_p(r)\) between that word and \(\Phi\).
The strengthened hypothesis makes the resulting function of
\((s,r)\) smooth.  Its derivative in \(r\) at zero is smooth in
\(s\).  At each fixed \(s\), bounded linearity of the outer
\(\rho\)-operator identifies this derivative with the outer
operator applied to \(d_p\Phi(s)\).
Induction on \(l\) proves the strengthened property for every
ordered derivative word.  Applying it to a translated constant
vector gives (SPR.18), and specializing the parameters gives the
two stability assertions.

No interchange of two Lie derivatives is used.  This argument
also proves the literal
\texttt{smooth\_\allowbreak along\_\allowbreak curve\_\allowbreak products\_\allowbreak stable} for arbitrary maps
\(e_p\) satisfying its product-smoothness premise: the group and
one-parameter laws are not used in this particular stability proof.
% END SPR9

% BEGIN SPR10
\paragraph{SPR10: Hilbert derivatives transfer to supplied continuous representatives.}
\label{claim:SPR10}
Here isometry of \(\rho\) is not required.  Retain strong
orbit continuity, bounded linearity and the other representative
hypotheses of the chapter.
Assume \(\mu\) is right invariant, let \(c:\mathbf R\to G\)
be a continuous one-parameter subgroup, and suppose
\[
 \frac{d}{dt}\bigg|_{0}\rho(c(t))v=v',\qquad
 \operatorname{Rep}(v,\phi),\quad
 \operatorname{Rep}(v',\phi'),
\]
with \(\phi,\phi'\) continuous.
Then, for every \(x\in G\),
\[
          \frac{d}{dt}\bigg|_0\phi(xc(t))=\phi'(x).
 \tag{SPR.19}
\]

\emph{Proof.}
The one-parameter law transports the Hilbert derivative to
\(\rho(c(t))v'\) at every \(t\).  Hence
\[
 \rho(c(b))v-\rho(c(a))v
                  =\int_a^b\rho(c(t))v'\,dt.
\]
Right invariance, (SPR.2), and the given representative identities
give continuous representatives \(\phi(\,\cdot\,c(t))\) and
\(\phi'(\,\cdot\,c(t))\) for the translated vectors.
SPR2 identifies the continuous representative of the integral.
By SPR1's uniqueness, for every \(x\) and \(a\le b\),
\[
 \phi(xc(b))-\phi(xc(a))
                    =\int_a^b\phi'(xc(t))\,dt.
 \tag{SPR.20}
\]
Using reversed endpoints covers \(b<a\).
Thus
\(\phi(xc(t))=\phi(x)+\int_0^t\phi'(xc(r))\,dr\).
The ordinary fundamental theorem, whose integrand is continuous,
gives (SPR.19).  This proves the full
\texttt{SmoothRep.\allowbreak hasDerivAt\_\allowbreak rep}; it does not differentiate a
chosen almost-everywhere representative without first constructing
the continuous representative of its derivative.
% END SPR10

% BEGIN SPR11
\paragraph{SPR11: one continuous representative assignment from finitely many averages.}
\label{claim:SPR11}
Assume the hypotheses of SPR6.  Suppose all admissible \(e_p\)
are continuous one-parameter subgroups commuting with every
\(h\in H\).  Let
\[
 \mathcal S=\{v:\operatorname{Sm}(v),\ \rho(h)v=v\ (h\in H)\}.
\]
Choose an admissible list \(q_1,\ldots,q_N\), \(s>0\), and put
\(\Theta(t)=e_{q_1}(t_1)\cdots e_{q_N}(t_N)\).
Assume its cube \(Q=[0,s]^N\) satisfies (SPR.11).
Then there is one assignment
\[
                     R:\mathcal H\longrightarrow\mathbf C^G
\]
such that \(Rv\) is continuous and
\(\operatorname{Rep}(v,Rv)\) for every \(v\in\mathcal S\).

\emph{Proof.}
SPR9 proves smoothness of \(d_pv\).  Its \(H\)-fixedness follows by
differentiating
\[
       \rho(h)\rho(e_p(r))v=\rho(e_p(r))\rho(h)v
                           =\rho(e_p(r))v.
\]
The operator \(\rho(h)\) is bounded, so it passes through this
Hilbert derivative.  Thus \(\mathcal S\) is stable under every
admissible \(d_p\) and every ordered word.

Use SPR8 for the fixed cube.  For each of its finitely many
weights \(\kappa_I\), SPR6 supplies evaluations
\(\Lambda_{I,x}\).  Define, for \(v\in\mathcal S\),
\[
                       Rv(x)=\sum_I\Lambda_{I,x}(D_Iv).
 \tag{SPR.21}
\]
Extend this assignment arbitrarily, for example by zero, outside
\(\mathcal S\).  Each \(D_Iv\) is fixed by \(H\), so every term
of (SPR.21) is continuous.
SPR6 identifies its representative as the corresponding averaged
Hilbert vector.  Summing and applying (SPR.16) gives
\(\operatorname{Rep}(v,Rv)\).

The same finite weights, derivative words, and evaluation
functionals are used for every \(v\in\mathcal S\).
Consequently derivatives of different vectors do not require
unrelated pointwise choices.
% END SPR11

% BEGIN SPR12
\paragraph{SPR12: the same representative has every original pointwise derivative.}
\label{claim:SPR12}
For the single assignment \(R\) in SPR11, every
\(v\in\mathcal S\) satisfies
\[
\begin{split}
 &t\longmapsto Rv\bigl(g\,e_{p_1}(t_1)\cdots e_{p_k}(t_k)\bigr)
                                  \text{ is }C^\infty,\\
 &R_{p_1}\cdots R_{p_m}(Rv)=R(d_{p_1}\cdots d_{p_m}v),
\end{split}
\tag{SPR.22}
\]
for every admissible list and every \(g\in G\), where
\(R_p\phi(x)=\frac{d}{dt}|_0\phi(xe_p(t))\).
In the second line both sides are continuous and represent the
same Hilbert derivative word.

\emph{Proof.}
Put \(a(t)=e_{p_1}(t_1)\cdots e_{p_k}(t_k)\).
SPR9 and commutation with \(H\) show that \(\rho(a(t))v\)
belongs to \(\mathcal S\).
The representative translation identity and uniqueness imply
\[
                     Rv(ga(t))=R(\rho(a(t))v)(g).
 \tag{SPR.23}
\]
Using (SPR.21), the right-hand side is
\[
             \sum_I\Lambda_{I,g}\bigl(D_I(\rho(a(t))v)\bigr).
\]
Every vector-valued summand is smooth by SPR9, and
\(\Lambda_{I,g}\) is bounded linear.  This proves the first line
of (SPR.22).  It avoids a continuity assertion for an unbounded
derived operator on the ambient Hilbert space.

For a single derivative, SPR10 applies to the continuous
representatives \(Rv\) and \(R(d_pv)\), supplied simultaneously
by SPR11.  It gives \(R_p(Rv)=R(d_pv)\).
The set \(\mathcal S\) is derivative-stable, so induction on a
right-fold list gives the second line in its original order.
At each induction step the new representative is the value of
this same \(R\).  This also proves the source's
\texttt{foldr\_\allowbreak rel} installation and the complete generic
\texttt{exists\_\allowbreak smooth\_\allowbreak rep}.
% END SPR12

% BEGIN SPR13
\paragraph{SPR13: the full original smooth automorphic representative theorem.}
\label{claim:SPR13}
Fix the original \(D,\theta,V\), with \(\operatorname{IsInfFam}(\theta)\),
the same \(\operatorname{FinPair}(V,\mu,\mathcal F)\), and any
\(\operatorname{IsReg}(V,\mu,\mathcal F,\rho)\).
Fix one open subgroup \(U\le G=G_A(D,V)\).
If \(F\in L^2(\operatorname{AutQ}(V),\nu_Q)\) satisfies the original
\(\operatorname{IsSmoothVector}'\) and is fixed by
\(U\cap G_f(V)\), there is one \(\psi:G\to\mathbf C\) such that
\[
 \operatorname{Repr}(F,\psi),\qquad
                  \operatorname{IsAutFun}(V,\psi),
 \tag{SPR.24}
\]
all original finite products of \(\operatorname{oneP}\)-curves give
smooth functions after every left base point, and, for every
admissible ordered list \(l\),
\[
       R_l\psi\text{ is continuous},\qquad
                    \operatorname{Repr}(d_lF,R_l\psi).
 \tag{SPR.25}
\]
The \(R_l\) and \(d_l\) here are exactly the source's right folds.
This is the complete original
\texttt{SmoothRep.\allowbreak exists\_\allowbreak smooth\_\allowbreak repr}.

\emph{Proof.}
RQC1 supplies the group topology; RQC3 and RQC6 supply the literal
quotient measure, global representative convention, and any original
regular action.  Choose measurable representatives of the quotient
\(L^2\) classes and pull them back by \(\operatorname{autMk}\):
\[
                   \iota(F)(g)=F(\operatorname{autMk}(g)).
\]
They are exactly left rational invariant.  A different quotient
representative changes the pullback only on a globally Haar-null
set by RQC3, so this assignment gives precisely the source's
\(\operatorname{Repr}\), including for its chosen \(L^2\) functions.
Linearity almost everywhere and (SPR.2) follow from HCU16 and
\(\operatorname{IsReg.val}\).

RQC11 gives (SPR.1) for every compact \(K\), with
\[
       M_K=\#(G_Q(V)\cap KK^{-1})\,\mu(\mathcal F)^{1/2}.
 \tag{SPR.26}
\]
This step needs no quantitative rational-point growth estimate.
It applies to the measurable rational-invariant pullbacks just
chosen, as in the full scope of RQC11.

Set \(H=U\cap G_f(V)\).
The finite subgroup is the kernel of the archimedean homomorphism,
hence is normal.  Therefore \(y^{-1}Hy\subseteq H\) for \(y\in U\).
The original real one-place elements commute with \(G_f(V)\),
by the coordinate separation in HCU14.
Take \(P\) to be the pairs \((\tau,X)\) with
\(X\in\operatorname{lieAt}(V,\tau)\), and
\(e_p(t)=\operatorname{oneP}(\theta,V,\tau,X,t)\).
ASB3 gives their one-parameter laws and the actual finite ordered
chart; ASB4 gives its exact scalar comparison (SPR.11).
All hypotheses of SPR11--SPR12 are now checked on the original
objects.

Let \(\psi=RF\).  Equations (SPR.24)--(SPR.25), except for the
finite and rational invariances in \(\operatorname{IsAutFun}\),
follow immediately from SPR11--SPR12.
For a rational \(\gamma\), the continuous functions
\(\psi(\gamma\,\cdot)\) and \(\psi\) represent \(F\), by left
Haar invariance and
\(\operatorname{autMk}(\gamma g)=\operatorname{autMk}(g)\).
They are equal everywhere by SPR1.
For \(k\in U\cap G_f(V)\), the continuous functions
\(\psi(\,\cdot\,k)\) and \(\psi\) likewise represent the same
\(F=\rho(k)F\), so they are equal everywhere.
Thus the original open \(U\) witnesses the finite level; it has
not been replaced by a level depending on \(l\).
The smooth curve statement is exactly (SPR.22), and its derivative
identity gives the whole source list quantifier in (SPR.25).

This proof uses the source's fixed finite pair and global
almost-everywhere convention.  It proves existence of the
representatives and every derivative identity, not existence
of the finite pair itself.
% END SPR13

\subsection{Installation in the original multiplicity space}

For SPR14--SPR21 fix the original \(D,\theta,V,S,M,w,v\)
with \(\operatorname{IsInfFam}(\theta)\),
\(\operatorname{HolShape}(M)\), \(\operatorname{WtIs}(M,w)\), and
\(\operatorname{IsExtVec}(M,v)\).
Fix one \(\operatorname{FinPair}(V,\mu,\mathcal F)\), put
\(\mathcal H=L^2(\operatorname{AutQ}(V),\nu_Q)\), and take the
original regular action supplied by RQC6.
Let \(J\subseteq G_f(V)\), and retain
\[
 \operatorname{HolForm}(\theta,M,w,J,\chi,\phi),
                 \qquad \operatorname{ZCusp}(V,\phi).
 \tag{SPR.27}
\]
Thus \(\phi\) is continuous, rational-left invariant, bounded,
of the prescribed finite character and original holomorphic type.
The separate \(\operatorname{ZCusp}\) hypothesis is retained.
HCU16--HCU18 give the original \(F_0\) and bounded multilinear map
\[
       \Phi:\prod_\tau M.\mathcal H_\tau\longrightarrow\mathcal H,
       \qquad \Phi(v)=F_0,\quad
                     \operatorname{Repr}(F_0,\phi),
 \tag{SPR.28}
\]
with all its original equivariance, closed-orbit and finite-character
properties.
Let \(U\) be one open subgroup witnessing the finite invariance
of \(\phi\).  Every \(\Phi(y)\) is fixed by \(U\cap G_f(V)\)
by HCU18.  RQC9--RQC10 give its Hilbert smoothness and all derivative
words.  The same \(U\) will be used throughout this installation.

% BEGIN SPR14
\paragraph{SPR14: the pointwise realization is one multilinear map.}
\label{claim:SPR14}
For \(y=(y_\tau)_\tau\in M.\mathrm{Vec}=\prod_\tau M.\mathrm{sm}_\tau\),
let \(T(y)\) be the continuous representative of
\(\Phi((y_\tau)_\tau)\) supplied by SPR13.
These choices form one complex multilinear map
\[
                    T:M.\mathrm{Vec}\longrightarrow\mathbf C^G,
 \tag{SPR.29}
\]
and every \(T(y)\) satisfies the original
\(\operatorname{IsAutFun}\), at the same \(U\), and all smooth
curve-product conditions.

\emph{Proof.}
Existence follows from RQC9, the common finite fixedness in
HCU18, and SPR13.  The construction may use the single \(R\)
of SPR11 installed in SPR13; alternatively, uniqueness makes
any choices agree.
Fix a place \(\tau\), the other entries of \(y\), and two smooth
vectors \(a,b\) in that place.
The continuous functions
\(T(y[\tau\leftarrow a+b])\) and
\(T(y[\tau\leftarrow a])+T(y[\tau\leftarrow b])\)
represent the same Hilbert vector by multilinearity of \(\Phi\)
and HCU16.  They are equal everywhere.
The scalar identity follows in the same way.
These are every defining field of the source's
\(\operatorname{MMap}(M)\).  SPR13 supplies the asserted
automorphic and smooth properties for each value.
% END SPR14

% BEGIN SPR15
\paragraph{SPR15: exact finite character and original one-place equivariance.}
\label{claim:SPR15}
The map \(T\) satisfies
\[
 T(y)(xk)=\chi(k)T(y)(x)\quad(k\in J)
 \tag{SPR.30}
\]
and, whenever \(g\in U(\operatorname{gramAt}(V,\tau))\) and
\(z\in M.\mathrm{sm}_\tau\) satisfies
\(z=M.\mathrm{rep}_\tau(g)y_\tau\),
\[
             T(y[\tau\leftarrow z])
              =r\operatorname{Tr}\bigl(\operatorname{atTau}
                    (\theta,V,\tau,g)\bigr)T(y).
 \tag{SPR.31}
\]
These are precisely the \texttt{lev} and \texttt{equiv} fields
of the original \(\operatorname{IsMultMap}\).

\emph{Proof.}
The finite-character clause of HCU18 gives
\(\rho(k)\Phi(y)=\chi(k)\Phi(y)\), since the same identity holds
for \(\phi\) and \(k\in J\subseteq G_f(V)\).
HCU16 identifies the continuous representatives of these two
Hilbert vectors as the two sides of (SPR.30), so continuous
uniqueness gives the identity at every \(x\).
No continuity or multiplicativity of the function \(\chi\)
outside the given source hypotheses is inserted.

For (SPR.31), HCU18 gives the exact one-place identity for
\(\Phi\) with the original \(\operatorname{atTau}\) and frame.
Both sides of (SPR.31) are continuous by SPR14 and continuity
of right translation.  HCU16 identifies both with that same
Hilbert vector.  Uniqueness proves (SPR.31).
This uses the actual updated tuple of smooth vectors, as
quantified in the source; it asserts no extension of \(T\) to
all ambient Hilbert vectors.
% END SPR15

% BEGIN SPR16
\paragraph{SPR16: pointwise differentiation of the same map in every ordered word.}
\label{claim:SPR16}
For \(X\in\operatorname{lieAt}(V,\tau)\),
\[
 R_{\tau,X}T(y)
   =T\bigl(y[\tau\leftarrow d\operatorname{Rep}_\tau(X)y_\tau]\bigr).
 \tag{SPR.32}
\]
For every admissible list \(l\), let \(l_\tau\) be its subsequence
at \(\tau\), retaining its order, and put
\[
        y^l_\tau=d\operatorname{Rep}_{l_\tau}y_\tau.
\]
Then \(y^l\in M.\mathrm{Vec}\) and
\[
                 R_lT(y)=T(y^l),\qquad
                   d_l\Phi(y)=\Phi(y^l).
 \tag{SPR.33}
\]

\emph{Proof.}
RQC7 gives stability of the original \(M.\mathrm{sm}_\tau\)
under every allowed derived action.
RQC9 identifies the Hilbert derivative
\(d_{\tau,X}\Phi(y)\) with the corresponding updated \(\Phi\)-value.
SPR13 says that \(R_{\tau,X}T(y)\) is continuous and represents
that Hilbert derivative.
The updated \(T\)-value has the same properties by SPR14.
HCU16's uniqueness gives (SPR.32).

Induct on the original right-fold list.
At its first letter update the corresponding factor of the tuple
already produced by the rest of the list.
The filtered subsequence within each place retains precisely this
order.  This is the existing RQC10 word formula for the Hilbert
map, now with its pointwise counterpart justified by (SPR.32).
No same-place derivatives have been permuted.
The same finite-level fixedness continues to hold for every
\(\Phi(y^l)\), by HCU18, independently of \(l\).
% END SPR16

% BEGIN SPR17
\paragraph{SPR17: one growth exponent for every value and every derivative word.}
\label{claim:SPR17}
For the fixed original finite level \(U\), there is one natural
number \(N\) such that for every \(y\in M.\mathrm{Vec}\), every
admissible word \(l\), and every \(x\in G\), there is a constant
\(C_{y,l,x}\) with
\[
              |R_lT(y)(xu)|
                 \le C_{y,l,x}\operatorname{normInf}(u)^N
                       \qquad(u\in G_\infty(V)).
 \tag{SPR.34}
\]
The height is the literal original sum of the archimedean matrix
and inverse-matrix entry norms.

\emph{Proof.}
ASB14, applied once to this quotient and this open \(U\), supplies
one \(N\), one finite list \(L\) of admissible derivative words,
and, for every \(x\), a constant \(C_x\) before the Hilbert vector
and its continuous representative are varied:
\[
 |\psi(xu)|
 \le C_x\operatorname{normInf}(u)^N
                           \sum_{a\in L}\|d_aF\|_2.
 \tag{SPR.35}
\]
For a fixed outer word \(l\), take
\(F=\Phi(y^l)\) and \(\psi=T(y^l)=R_lT(y)\).
Their smoothness, same-level fixedness and representative identity
are supplied by RQC9, HCU18 and SPR14--SPR16.
Thus (SPR.35) proves (SPR.34), for example with
\[
                  C_{y,l,x}
                  =|C_x|\sum_{a\in L}\|d_a\Phi(y^l)\|_2.
\]
The finitely many displayed norms are finite Hilbert norms;
no uniform bound in \(l\) for them is required by the source.

This supplies the full \texttt{growth} field of
\(\operatorname{IsAutFunSm}\), with \(N\) chosen before both
\(y\) and \(l\).
Its uniformity is for this fixed finite level.  If the initial
seed or external parameter changes the level, this argument
makes no additional assertion of a common \(N\) across those
levels.  The earlier TAE kernel exponent retains its separate
stronger parameter order.
% END SPR17

% BEGIN SPR18
\paragraph{SPR18: the central-finiteness field on the actual pointwise realization.}
\label{claim:SPR18}
For every \(y\), the finite-dimensional space
\[
                         W_y=\mathbf C\,T(y)
 \tag{SPR.36}
\]
contains \(T(y)\) and is stable under every original
\(\operatorname{gelf}(\theta,V,\tau,k)\).
Thus the source's full \texttt{zfin} field holds.

\emph{Proof.}
The original \(\operatorname{HolShape}\),
\(\operatorname{WtIs}\) and \(\operatorname{IsExtVec}\)
give the exact extremal Borel equations by HCU12.
RQC20 therefore supplies, for each \(\tau,k\), one scalar
\(c_{\tau,k}\) for the corresponding \(k+1\)-cycle operator on
every smooth model vector.
The now established identity (SPR.32) checks the previously
retained real pointwise-derivative premise of RQC21.
Its complexification and ordered-word transfer give
\[
             \operatorname{gelf}(\theta,V,\tau,k)T(y)
                          =c_{\tau,k}T(y).
\]
The same formula holds on scalar multiples of \(T(y)\), by
the smooth derivative identities, so \(W_y\) is stable.
Its dimension is at most one, including when \(T(y)=0\).
The operator normalization and right-fold order are exactly
those already verified in RQC17--RQC21.
This use of central finiteness occurs after the bounded seed
and Hilbert map have been supplied.
% END SPR18

% BEGIN SPR19
\paragraph{SPR19: full cuspidality of every value of the same map.}
\label{claim:SPR19}
For every \(y\in M.\mathrm{Vec}\),
\[
             \operatorname{ZCusp}(V,T(y)),
                       \qquad \operatorname{IsCuspFun}(V,T(y)).
 \tag{SPR.37}
\]

\emph{Proof.}
HCU18 places \(\Phi(y)\) in the closed span of the full regular
orbit of \(F_0\).
The seed \(\phi\) is an original automorphic function with the
separate \(\operatorname{ZCusp}\) condition in (SPR.27).
SPR14 supplies the continuous rational-invariant representative
\(T(y)\) of \(\Phi(y)\), with its original finite-level law.
All hypotheses of the complete RQC15 cusp-closure statement
therefore hold on the same quotient and measure.
It proves every transvection, every Haar-measure and every
fundamental-domain quantifier in
\(\operatorname{ZCusp}(V,T(y))\).
CTC12, as installed in RQC16, gives genuine cuspidality.
This completes the \texttt{cusp} field; no restriction of a
bare \(L^2\) class to a unipotent subgroup occurs.
% END SPR19

% BEGIN SPR20
\paragraph{SPR20: the exact separate Sobolev continuity of point evaluation.}
\label{claim:SPR20}
For every place \(\tau\), tuple \(y\in M.\mathrm{Vec}\), and
\(x\in G\), there is an algebraic complex linear functional
\[
                   \ell:M.\mathcal H_\tau\longrightarrow\mathbf C
\]
satisfying the literal \(\operatorname{Sobolev}_2\) predicate and
\[
                    \ell(w)=T(y[\tau\leftarrow w])(x)
                        \qquad(w\in M.\mathrm{sm}_\tau).
 \tag{SPR.38}
\]
This is the full \texttt{cont} field of
\(\operatorname{IsMultMap}\).

\emph{Proof.}
Evaluate the ASB14 estimate at the identity real element.
For the same finite list \(L\) as in (SPR.35), there is \(B_x\ge0\)
such that
\[
             |T(z)(x)|\le B_x
                         \sum_{a\in L}\|d_a\Phi(z)\|_2
                           \qquad(z\in M.\mathrm{Vec}).
 \tag{SPR.39}
\]
One may take \(B_x=|C_x|\operatorname{normInf}(1)^N\).
No division by this height is used.

Fix the entries of \(y\) outside \(\tau\).
By RQC10 and bounded multilinearity of \(\Phi\), for each \(a\in L\),
\[
\begin{split}
 \|d_a\Phi(y[\tau\leftarrow w])\|_2
 &\le b_a\,\|d\operatorname{Rep}_{a_\tau}w\|,\\
 b_a&=\|\Phi\|\prod_{\sigma\ne\tau}
                    \|d\operatorname{Rep}_{a_\sigma}y_\sigma\|.
\end{split}
\tag{SPR.40}
\]
This uses the actual filtered word \(a_\sigma\), with its order
unchanged.
Let \(b\) be the maximum of zero and these finitely many \(b_a\)'s;
if \(L\) is empty let \(b=0\).
Equations (SPR.39)--(SPR.40) give
\[
 |T(y[\tau\leftarrow w])(x)|
 \le B_xb\sum_{a\in L}
                         \|d\operatorname{Rep}_{a_\tau}w\|.
 \tag{SPR.41}
\]
The list \([a_\tau:a\in L]\) is finite, has only admissible Lie
letters, and may retain repetitions, as permitted by the source.

By SPR14, evaluation in this slot is a linear functional on
\(M.\mathrm{sm}_\tau\).
Extend a vector-space basis of this subspace to a basis of
\(M.\mathcal H_\tau\), and extend the functional by zero on the
new basis vectors.
Equation (SPR.41) holds on every original smooth vector, since
such a vector belongs to \(M.\mathrm{sm}_\tau\).
It is precisely the \(\operatorname{Sobolev}_2\) bound for this
algebraic extension \(\ell\).
No boundedness for the Hilbert norm alone is asserted or needed.
This proves \texttt{BackForm.\allowbreak cont\_\allowbreak value}, including its full
ambient algebraic-functional conclusion.
% END SPR20

% BEGIN SPR21
\paragraph{SPR21: the complete original bounded-holomorphic multiplicity producer.}
\label{claim:SPR21}
For the original data, assume
\[
\begin{gathered}
 \operatorname{IsInfFam}(\theta),\
 \operatorname{HolShape}(M),\
 \operatorname{WtIs}(M,w),\
 \operatorname{IsExtVec}(M,v),\\
 \exists\mu,\mathcal F\ \operatorname{FinPair}(V,\mu,\mathcal F),\\
 J\subseteq G_f(V),\
 \operatorname{HolForm}(\theta,M,w,J,\chi,\phi),\
 \operatorname{ZCusp}(V,\phi).
\end{gathered}
\tag{SPR.42}
\]
Then there is an original \(T:\operatorname{MMap}(M)\) with
\[
                   \operatorname{IsMultMap}(\theta,M,J,\chi,T),
                                  \qquad T(v)=\phi.
 \tag{SPR.43}
\]
This is the full source theorem
\texttt{BackForm.\allowbreak exists\_\allowbreak multMap}, without a remaining unnamed
smooth-representative, evaluation or Sobolev premise.

\emph{Proof.}
Choose one finite pair from (SPR.42), once for the whole map.
RQC6 gives its original regular action.
HCU16 turns the bounded seed into \(F_0\);
HCU17 checks its exact extremal coefficient identity;
HCU18 constructs the same-quotient \(\Phi\) in (SPR.28).
Choose one finite level from the seed's
\(\operatorname{IsAutFun}\) and use SPR14 to construct \(T\).

Its \texttt{aut} field consists of the automorphic and smooth
properties in SPR14, central finiteness in SPR18, and the
full all-word growth statement in SPR17.
SPR19 gives \texttt{cusp}.
SPR15 gives \texttt{lev} and \texttt{equiv}.
SPR20 gives \texttt{cont}.
These are all five fields of the literal source
\(\operatorname{IsMultMap}\).

Finally \(\Phi(v)=F_0\), and both \(T(v)\) and \(\phi\) are
continuous representatives of \(F_0\).
HCU16 gives \(T(v)=\phi\) at every point.
The zero seed causes no exception: the same construction gives
the zero map.
The boundedness in \(\operatorname{HolForm}\), the initial
finite pair, and the seed \(\operatorname{ZCusp}\) were genuine
inputs.  None has been deduced from the downstream
central-finiteness conclusion.
% END SPR21

% BEGIN SPR22
\paragraph{SPR22: equality at one nonzero lowest-type tuple gives full restriction.}
\label{claim:SPR22}
Let the original model satisfy \(\operatorname{LowSm}(M)\) and
\(\operatorname{LowIrr}(M)\).  Let \(T\) satisfy
\(\operatorname{IsMultMap}(\theta,M,J,\chi,T)\), and let
\(U:\operatorname{KMap}(M)\) satisfy the original compact
equivariance \(\operatorname{KEquiv}(\theta,M,U)\).
If every component of a lowest-type tuple \(e\) is nonzero and
\[
       T((e_\tau)_\tau)=U(e),
 \tag{SPR.44}
\]
then \(\operatorname{RestrictsTo}(M,T,U)\): equality holds for
every lowest-type tuple, with the original embeddings into
\(M.\mathrm{sm}\).

\emph{Proof.}
For each \(\tau\), the linear span of the compact-group orbit of
\(e_\tau\) is nonzero and invariant.  Irreducibility makes it
the entire original lowest-type space.  This is an algebraic
span; every vector is a finite linear combination of such orbit
vectors.  Put \(B=T|_{\mathrm{low}}-U\), a multilinear map to
the space of functions on the original adelic group.

We first check that its zero relation propagates under the
actual compact action.  Given a compact matrix \(k\) at \(\tau\),
put \(g=(M.\mathrm{fr}_\tau)^{-1}kM.\mathrm{fr}_\tau\).
The source identity
\[
 \mathrm{gramAt}_\tau
   =(M.\mathrm{fr}_\tau)^*J_{M.\varepsilon_\tau}
                         M.\mathrm{fr}_\tau
 \tag{SPR.45}
\]
and \(k^*J_{M.\varepsilon_\tau}k=J_{M.\varepsilon_\tau}\)
give \(g^*\mathrm{gramAt}_\tau g=\mathrm{gramAt}_\tau\).
The definition of the model representation gives
\(M.\operatorname{rep}_\tau(g)=M.\sigma_\tau(k)\).
Thus the \texttt{equiv} field of \(T\) and compact equivariance
of \(U\) give the same right translate after updating this
slot.  In particular, if \(B(v)=0\), then
\(B(v[\tau\leftarrow M.\sigma_\tau(k)v_\tau])=0\).
The embeddings of an updated tuple agree slot by slot, so this
is an identity of the actual multilinear maps.

Starting from (SPR.44), propagate through compact translates
in each of the finitely many slots.  Expand an arbitrary tuple
as finite linear combinations of those translates in each slot.
Multilinearity makes every term zero.  Hence \(B=0\).
This proves the full source \texttt{ThetaMap.\allowbreak restrictsTo\_\allowbreak of\_\allowbreak eq};
there is no additional density, completion, or continuity step.
% END SPR22

% BEGIN SPR23
\paragraph{SPR23: the lowest-type extension after its two initial analytic inputs.}
\label{claim:SPR23}
Retain \(\operatorname{IsInfFam}(\theta)\),
\(\operatorname{HolShape}(M)\), \(\operatorname{WtIs}(M,w)\),
\(\operatorname{LowSm}(M)\), \(\operatorname{LowIrr}(M)\),
and one actual extremal tuple \(v\) with
\(\operatorname{IsExtVec}(M,v)\).
Let \(e\) be its tuple in the original lowest types.
Let \(J\subseteq G_f\), and let \(U:\operatorname{KMap}(M)\)
satisfy, for every lowest-type tuple \(z\),
\[
\begin{gathered}
 \operatorname{IsAutFun}(V,U(z)),\qquad
 U(z)(gk)=\chi(k)U(z)(g)\quad(k\in J),\\
 \operatorname{IsSmoothFn}(V,\theta,U(z)),\qquad
 \operatorname{KEquiv}(\theta,M,U),\\
 R^\mathbb C_{\tau,Z}U(z)=0
                  \quad(Z\in\mathfrak p^-_{M,\tau}).
\end{gathered}
\tag{SPR.46}
\]
Assume the seed \(U(e)\) is \(\operatorname{ZCusp}\), and supply
the following two initial analytic outputs explicitly:
\[
 \exists\mu,\mathcal F\ \operatorname{FinPair}(V,\mu,\mathcal F),
 \qquad \exists C\ \forall g\ |U(e)(g)|\le C.
 \tag{SPR.47}
\]
Then there is \(T:\operatorname{MMap}(M)\) satisfying
\(\operatorname{IsMultMap}(\theta,M,J,\chi,T)\) and
\(\operatorname{RestrictsTo}(M,T,U)\).

\emph{Proof.}
HCU19 checks the original \(\operatorname{HolAd}\) equations
for \(U(e)\), from the matching weight, compact action,
smoothness and lowering equations in (SPR.46).
Together with \(\operatorname{IsAutFun}\), the level law and
the bound in (SPR.47), these are all fields of
\(\operatorname{HolForm}\).  Apply SPR21 with the same finite
pair and the separately supplied \(\operatorname{ZCusp}\).
It gives \(T(v)=U(e)\).  The unit-norm clause in
\(\operatorname{IsExtVec}\) implies \(e_\tau\ne0\) for every
\(\tau\).  SPR22 gives the full restriction identity.

This reconstructs the complete extension consumer after
(SPR.47).  The original \texttt{ThetaMap.\allowbreak exists\_\allowbreak extension}
obtains those two outputs from target quotient volume and
boundedness of smooth cuspidal uniformly moderate forms.
Neither supplier is asserted here.  In particular, an
\(\operatorname{IsHolCuspKMap}\) packet does not silently supply
the open finite level in \(\operatorname{IsAutFun}\); that field
is explicit in (SPR.46).
% END SPR23

% BEGIN SPR24
\paragraph{SPR24: full linear choice from pointwise theta firing.}
\label{claim:SPR24}
Let \(X,M,N\) be complex vector spaces, \(W\le M\) a subspace,
and \(r:M\to N\), \(u:X\to N\) linear maps.  Then
\[
 (\forall x\ \exists m\in W\quad r(m)=u(x))
 \quad\Longrightarrow\quad
 \exists y:X\to M\text{ linear},\quad
 y(X)\subseteq W,\quad r\circ y=u.
 \tag{SPR.48}
\]
Consequently the full original
\texttt{ThetaFiresLinear.\allowbreak exists\_\allowbreak linear\_\allowbreak lift} holds, with its
bilinear theta integral on \(\operatorname{range}(\mathrm{ext})\),
for an arbitrary source vector space \(X\).  No finite-dimensional
or injectivity premise is needed.

\emph{Proof.}
Choose a Hamel basis \((b_j)\) of \(X\).  For each basis vector
choose \(m_j\in W\) with \(r(m_j)=u(b_j)\).
For the unique finite basis expansion \(x=\sum_jc_jb_j\), set
\(y(x)=\sum_jc_jm_j\).  This is linear, lies in \(W\), and
has the required image under \(r\).  This proves (SPR.48).

We verify every input when using this choice for the actual
theta integral.  Let \(\nu,\mathcal F\) be its unchanged
measure and fundamental-domain set.  Write its kernel as
\(\mathrm{th}(g,h,s)=\mathrm{thL}(g,h)(s)\), with
\(\mathrm{thL}(g,h)\) complex linear, and its polynomial map as
the actual linear \(\mathrm{ptL}\).
Let \(\mathrm{io}:TK\to\operatorname{Hom}_{\mathbb C}
(\mathrm{Tau},S_\infty)\) and
\(\mathrm{ext}:X\to(H'\to\mathrm{Tau})\) be the original
linear maps.  Assume the integrands are integrable on
\(\mathcal F\) for every \(x,e,g\).
For \(f\in\operatorname{range}(\mathrm{ext})\), put
\[
 P(f)(e)(g)=\int_{\mathcal F}
   \mathrm{thL}(g,h)
     \bigl(\mathrm{ptL}(\mathrm{io}(e)(f(h)))\bigr)\,d\nu(h).
 \tag{SPR.49}
\]
Such an \(f\) has a preimage, so its integrability follows from
the hypothesis.  Formula (SPR.49) uses \(f\) itself; it does not
choose a preimage in its definition.  Additivity of integrals
of integrable functions and their scalar law prove that \(P\)
is linear in \(f\) and in \(e\).  It is exactly the original
\(\operatorname{Lift.IsThetaP}\) map, including the integrability
field and the pointwise integral equality.

Let \(M\) now be the space of multilinear maps from the model
smooth spaces to functions on the target group, and let
\(W\) be the original prescribed subspace of such maps.
Take \(N\) to be all functions of a lowest-type tuple with
values in target functions.  Restriction defines a linear
\(r:M\to N\).  The rule
\[
 u(x)(z)=P(\mathrm{ext}(x))
                       \bigl(\mathop{\otimes}_\tau z_\tau\bigr)
 \tag{SPR.50}
\]
is linear in \(x\).  Pointwise firing supplies, for each \(x\),
an \(m\in W\) whose restriction is (SPR.50).  Apply (SPR.48).
The resulting \(y\), together with this same \(P\), has every
field of \texttt{exists\_\allowbreak linear\_\allowbreak lift}.
For the original theta kernel, \(\mathrm{thL}\) is evaluation
at zero after the original Weil operator, and
\(\mathrm{ptL}\) is the original polynomial-to-Schwartz map;
their linearity is literal.  Thus this also proves the complete
\texttt{ThetaFiresOfMult.\allowbreak exists\_\allowbreak isThetaMult\_\allowbreak of\_\allowbreak firesAt}
consumer, for every fixed original \(i,t\).
% END SPR24

% BEGIN SPR25
\paragraph{SPR25: the same-realization theta map with the exact two remaining analytic gates.}
\label{claim:SPR25}
Fix the original common \(R\), with its level, family and model
laws, and its actual \(\mathrm{thP\_\omega P}\),
\(\mathrm{thP\_vP}\), \(\mathrm{thP\_pt\_theta}\) and
\(\mathrm{thP\_ioP}\) laws.  The last two may be supplied by
JIO19 when its original weak stored-choice laws are available.
Choose the original unit extremal tuples \(e_i\) from HCU20,
before \(t\).  Denote by \(U_{i,t,x}\) the literal lowest-type
theta integral of TAE13, for \(x\in(R.\mathrm{HV}_{i,t}).H\).
Assume precisely the following additional outputs:
\[
 \forall i\ \exists\mu_i,\mathcal F_i\quad
       \operatorname{FinPair}(W_i,\mu_i,\mathcal F_i).
 \tag{SPR.51}
\]
\[
 \forall i,t,x\ \exists C_{i,t,x}\ \forall g\quad
       |U_{i,t,x}(e_i)(g)|\le C_{i,t,x}.
 \tag{SPR.52}
\]
Then there is one family of linear theta maps \(\mathrm{th}_{i,t}\)
into the original \(\operatorname{MMap}(R.\mathrm{MW}_i)\)
such that updating \(R.\mathrm{thF}\) with this family satisfies
both original \(\mathrm{thP\_thF\_mem}\) and
\(\mathrm{thP\_thF}\), for all \(i,t\).

\emph{Proof.}
HCU20 supplies the matching holomorphic shape, integral weight,
lowest-type smoothness and irreducibility, and the same unit
extremal tuple.  TAE13 supplies the actual \(U_{i,t,x}\), all
its smoothness and lowering equations, its compact equivariance,
and its uniform moderate-growth packet.  Its left invariance
and continuity are those of the original theta integral.
The principal-level law from BTH5 and BTH15 uses exactly
\(\chi=1\), as in the original source.

For clarity the open-level input is checked separately.
Write \(L\) for the original lattice-level datum, \(\beta\) for
its actual conjugating matrix, and \(T=L.P_g\cup L.Q\).
The integral finite tensor subring is open in the original module
topology: finite place coordinates identify it with local integer
rings, and any fixed integral basis gives the same open lattice
outside finitely many places; at those remaining places two full
lattices are commensurable and open.
Near \(k=1\), require the entries of both
\(m(k)=\beta^{-1}\mathrm{finE}(k)\beta\) and \(m(k)^{-1}\)
to be integral, and require
\[
  m(k)_{ab}-\delta_{ab}
       \in(\operatorname{Norm}P)\,\mathrm{intAf}
            \quad(P\in T,\ a,b).
\]
These finitely many conditions hold on an open identity
neighborhood, by continuity, inversion and multiplication by
the nonzero field scalar \((\operatorname{Norm}P)^{-1}\).
Integrality of both matrices puts \(\mathrm{finE}(k)\) in the
original lattice stabilizer.  Since \(\operatorname{Norm}P\)
annihilates its residue field, its reduction at every \(P\in T\)
is the identity.  It therefore satisfies both original block-flag
conditions and the original residual subgroup condition, because
each contains the identity.  This is the exact inclusion
\(K_g(\beta,T)\subseteq\mathrm{lev}(\beta,L)\).
The inverse image of the latter subgroup under the continuous
finite-component homomorphism contains this neighborhood, hence
is an open subgroup \(O\) of the full adelic group.
By the definition of the original \(J_i=\mathrm{levS}(\beta,L)\),
\(O\cap G_f\subseteq J_i\).  This proves the source
\texttt{ThetaMap.\allowbreak level\_\allowbreak open}, with its original conjugating
matrix and lattice datum.  Consequently every \(U_{i,t,x}(z)\)
has the required \(\operatorname{IsAutFun}\) field.  No new
level or character replaces \(J_i\).

The Heisenberg-averaging and constant-term argument of TAE13
and CTC13 supplies the seed \(\operatorname{ZCusp}\).
With (SPR.51)--(SPR.52), SPR23 therefore gives a genuine
\(\operatorname{IsMultMap}\) extending each \(U_{i,t,x}\).
It belongs to \(R.\mathrm{Wall}_i\), since that subspace is the
span of the original maps satisfying \(\operatorname{IsMultMap}\).
Integrability for every tensor input and target point is the
original definite-source integral theorem in TAE6 and TAE13.
These facts give the complete \(\operatorname{FiresAt}(R,i,t,x)\)
packet, for every \(x\), with the original kernel, measure,
fundamental domain, polynomial map and \(\mathrm{io}_i\).

Apply SPR24 for each \(i,t\).  It yields a single linear map in
\(x\), its range in the same \(R.\mathrm{Wall}_i\), and the
exact bilinear theta-integral law on
\(\operatorname{range}(R.\mathrm{extV}_{i,t})\).
Choose these maps as a family \(\mathrm{th}_{i,t}\).
The two conclusions are precisely the two fields of the updated
\(R\) in the source \texttt{ThetaMap.\allowbreak exists\_\allowbreak thetaMap}.
Only \(\mathrm{thF}\) has been replaced; all other data are the
same.  When an original stored \(\mathrm{OrZero}\) law is used,
this existence activates its predicate by the weak-choice law;
it does not identify the stored witness with the chosen one.

This is a conditional reconstruction of that final consumer.
The initial seed bound (SPR.52) and target finite pair (SPR.51)
remain unresolved original producers.  The smooth representatives,
their full derivative words, pointwise growth, the five
multiplicity-map fields, lowest-type restriction and linear
assembly have now been reconstructed after those two inputs.
% END SPR25

